\clearpage
\section{Metodología}

% ============================================================
\subsection{Técnicas, Métodos y Procedimientos}
% ============================================================

Para el desarrollo de la aplicación se adopta la metodología RAD (Rapid Application Development / Desarrollo Rápido de Aplicaciones), seleccionando y adaptando sus fases al contexto de un sistema de gestión nutricional con inteligencia artificial. RAD permite construir prototipos funcionales de forma iterativa e incremental con retroalimentación continua, reduciendo los tiempos de desarrollo mediante componentes
reutilizables y una participación activa del usuario final, lo que resulta adecuado para este proyecto dado que los requisitos se consolidan a partir de la interacción directa con el actor principal (el nutricionista). Se aplican modelos UML (diagramas de casos de uso, modelo de dominio, diagramas de secuencia y diagrama de clases) junto con el modelo de arquitectura C4 en las etapas de Diseño del Usuario y Construcción Rápida, lo que permite mantener una documentación estructurada y coherente con la implementación del sistema. Además, para la evaluación del prototipo se emplea un cuestionario Likert aplicado a expertos en nutrición, que permite medir de forma cuantitativa la percepción de
adecuación clínica, claridad en recomendaciones, facilidad de uso, confianza en las recomendaciones generadas y eficiencia temporal del sistema. Tras la aplicación del instrumento se analizan los resultados para calcular un índice global de satisfacción y evaluar la pertinencia del prototipo en el contexto clínico nutricional.

% ============================================================
\subsection{Fases del Proyecto}
% ============================================================

\subsubsection*{Objetivo Específico 1: Diseñar el prototipo NutriAPP}

\textbf{1. Planificación de Requisitos}
\begin{itemize}
  \item \textbf{Requisitos funcionales y no funcionales:} definición del alcance del  sistema, incluyendo los módulos de registro de pacientes, mediciones antropométricas, bioquímicas, historial clínico, hábitos dietéticos y asistente IA con pipeline RAG; así como las condiciones de rendimiento, seguridad y usabilidad.
  \item \textbf{Identificación de actores:} El nutricionista es actor principal del sistema y definición de los objetivos
  funcionales de cada módulo.
  \item \textbf{Identificación del catálogo alimentario:} planificación de la recopilación y organización
  de datos nutricionales de alimentos ecuatorianos.
  \item \textbf{Modelo conceptual:} diagrama UML que identifica las entidades
  principales del dominio (Paciente, Antropometría, HistorialClínico,
  HábitosDietéticos, MensajeChat, Alimento) y sus relaciones.
\end{itemize}

\textbf{2. Diseño}
\begin{itemize}
  \item \textbf{Diagramas de casos de uso:} mapeo de las interacciones del
  nutricionista con los módulos del sistema para validar el alcance definido.
  
  \item \textbf{Modelo de dominio y comportamiento:} diagrama de clases con
  atributos y cardinalidades, y diagramas de secuencia que ilustran flujos  principales, incluyendo el pipeline RAG (mensaje del nutricionista
  $\rightarrow$ backend $\rightarrow$ API Groq con Llama~3.3 $\rightarrow$
  respuesta al frontend).
  \item \textbf{Arquitectura C4:} documentación de la arquitectura en cuatro
  niveles (Contexto, Contenedores, Componentes, Código) describiendo la
  interacción entre el frontend React, el backend Express y la base de datos
  PostgreSQL; complementada con diagramas de despliegue y de paquetes.
\end{itemize}

\textbf{3. Desarrollo}

\textit{3.1 Construcción}
\begin{itemize}
  \item \textbf{Backend:} implementación de API REST con Express.js y TypeScript; módulos de gestión de pacientes, datos clínicos, hábitos dietéticos y asistente IA con pipeline RAG (extracción de texto con pdf-parse e inferencia mediante Groq API).
  \item \textbf{Frontend:} desarrollo de interfaz clínica con React y TailwindCSS; módulos de registro de pacientes, antropometría, historial clínico, hábitos dietéticos y chat con asistente nutricional.
  \item \textbf{Base de datos:} definición del esquema relacional y migraciones en PostgreSQL mediante Prisma ORM.
\end{itemize}

\textit{3.2 Pruebas}
\begin{itemize}
  \item \textbf{Pruebas de caja negra:} validación de los endpoints críticos del backend (registro de mediciones, generación de recomendaciones, pipeline RAG con documentos PDF).
\end{itemize}

\textit{3.3 Control de versiones}
\begin{itemize}
  \item Gestión del código fuente mediante GitHub, con ramas \texttt{main} y
  \texttt{develop}.
\end{itemize}

\textbf{1.4 Documentación de resultados}
\begin{itemize}
  \item Redacción de la documentación técnica del prototipo: descripción de la arquitectura implementada, decisiones de diseño, manual de instalación y guía de uso del sistema NutriAPP.
\end{itemize}

\subsubsection*{Objetivo Específico 2: Evaluar la satisfacción de los nutricionistas}

\textbf{4. Evaluación con Likert}

\textit{4.1 Diseño del cuestionario}
\begin{itemize}
  \item \textbf{Diseño del cuestionario Likert:} elaboración de un cuestionario con $\geq$10~ítems que evalúa adecuación clínica, claridad de las recomendaciones, facilidad de uso, confianza y eficiencia temporal, distribuido mediante Google Forms.
\end{itemize}

\textit{4.2 Aplicación del cuestionario}
\begin{itemize}
  \item \textbf{Población objetivo:} nutricionistas que interactúen con el prototipo en el entorno de pruebas mediante los perfiles genéricos establecidos.
  \item Aplicación del cuestionario a los nutricionistas participantes.
\end{itemize}

\textit{4.3 Análisis del cuestionario}
\begin{itemize}
  \item Tabulación de respuestas, cálculo del índice global de satisfacción y análisis estadístico.
\end{itemize}

\textit{4.4 Documentación de resultados}
\begin{itemize}
  \item Redacción del informe final con los resultados obtenidos, conclusiones y recomendaciones derivadas de la evaluación de satisfacción del prototipo NutriAPP.
\end{itemize}

% ============================================================
\subsection{Consideraciones éticas}
% ============================================================

La presente investigación involucra diferentes aspectos éticos, cada
uno con sus propias consideraciones:

\textbf{Datos del paciente:}
El prototipo procesa datos genéricos y anonimizados del paciente como: peso, talla, edad, sexo, nivel de actividad física y enfermedades o condiciones metabólicas
relevantes con el fin de generar recomendaciones nutricionales
personalizadas. En el contexto de esta investigación, dichos datos
corresponden a perfiles de prueba representativos de la población
ecuatoriana, construidos con fines experimentales y no necesariamente recopilados de pacientes reales. 

\textbf{Catálogo alimentario nutricional ecuatoriano:}
La base de conocimiento del sistema RAG se construye a partir de datos
nutricionales de alimentos ecuatorianos obtenidos de fuentes oficiales de
acceso público, como la Encuesta Nacional de Salud y Nutrición~(ENSANUT). Estos datos
son de naturaleza nutricional y no constituyen información personal
identificable; su uso con fines de investigación académica está permitido
bajo las condiciones de acceso abierto de dichas fuentes, respetando en toda
publicación derivada la atribución a los autores originales.

\textbf{Uso responsable de la inteligencia artificial:}
El sistema emplea el modelo de lenguaje de código abierto dentro de un pipeline de generación aumentada por recuperación~(RAG),
utilizado exclusivamente como herramienta de apoyo como sugerencias de los planes alimenticios para la revisión, aceptación del nutricionista. En la fase de prototipado, el modelo se accede mediante la
APIS, que proveen inferencias sobre modelos open-source sin retención
permanente de datos según su política de privacidad~\cite{gavai2025,arshad2025};
en una implementación bajo un entorno de producción con datos clínicos reales, el modelo de IA deberá
desplegarse localmente, eliminando cualquier transmisión de información a servicios externos de terceros y
preservando la confidencialidad del paciente. La decisión final sobre el
plan nutricional recae siempre en el nutricionista, preservando el rigor y
la responsabilidad médica. Los resultados generados no serán empleados para
diagnóstico clínico autónomo ni para ningún uso que comprometa los derechos
fundamentales o la dignidad de las personas. El alcance del proyecto finaliza
en la evaluación de la satisfacción percibida por el nutricionista frente a
las recomendaciones generadas.

\textbf{Privacidad en la evaluación de satisfacción:}
El cuestionario Likert aplicado a los expertos en nutrición participantes se
generará de forma anónima; no se recopilarán datos de identificación
personal de los profesionales encuestados. La participación será voluntaria e informada, y los resultados se utilizarán exclusivamente con fines para esta investigación académica, sin exposición pública de las respuestas.

% ============================================================
\subsection{Resumen de la metodología}
% ============================================================

En la Tabla~\ref{tab:metodologia} se presenta el resumen de la metodología del proyecto,
con las actividades, entregables, métodos y herramientas correspondientes a cada objetivo
específico.

\clearpage
\begin{landscape}
\begin{table}[h]
\centering
\caption{Metodología aplicada al proyecto}
\label{tab:metodologia}
\small
\begin{tabular}{|c|p{5.5cm}|p{3.8cm}|p{4cm}|p{4cm}|}
\hline
\textbf{Objetivo} & \textbf{Actividades} & \textbf{Entregables} &
\textbf{Métodos / Técnicas} & \textbf{Herramientas / Tecnologías} \\
\hline

\multirow{4}{*}{O.E.1}
& Planificación de Requisitos (Funcionales, No Funcionales, Identificación de actores, Modelo conceptual) e Identificación del catálogo alimentario ecuatoriano.
& Documento de requisitos, modelo conceptual UML y catálogo de alimentos ecuatorianos.
& Análisis de requisitos y modelado conceptual.
& Visual Studio Code, StarUML. \\ \cline{2-5}

& Diseño de la aplicación (Diagramas de casos de uso, Modelo de dominio y comportamiento, Arquitectura C4).
& Diagramas UML (casos de uso, dominio, secuencia) y arquitectura C4.
& Modelado UML y arquitectura C4.
& StarUML, draw.io. \\ \cline{2-5}

& Desarrollo del prototipo web (backend, frontend, BD) con pipeline RAG y catálogo de alimentos ecuatorianos.
& Prototipo web NutriAPP funcional con pipeline RAG integrado.
& Desarrollo iterativo de software (RAD).
& Node.js, Express.js, TypeScript, React, TailwindCSS, PostgreSQL, Prisma ORM, Groq API. \\ \cline{2-5}

& Documentación técnica: arquitectura de la aplicación.
& Documentación técnica del sistema NutriAPP.
& Documentación de software.
& \LaTeX, Visual Studio Code, GitHub. \\
\hline

\multirow{4}{*}{O.E.2}
& Diseño del cuestionario Likert (adecuación clínica, claridad, facilidad de uso, confianza y eficiencia temporal).
& Cuestionario Likert ($\geq$10 ítems) en Google Forms.
& Diseño de instrumentos de evaluación.
& Google Forms, Microsoft Excel. \\ \cline{2-5}

& Aplicación del cuestionario a nutricionistas.
& Cuestionario aplicado y respuestas recopiladas.
& Recolección de datos primarios.
& Prototipo NutriAPP, Google Forms. \\ \cline{2-5}

& Tabulación y análisis estadístico descriptivo para calcular el índice global de satisfacción.
& Informe de satisfacción con índice global de satisfacción.
& Análisis estadístico descriptivo.
& Google Forms, Microsoft Excel. \\ \cline{2-5}

& Redacción del informe final con resultados, conclusiones y recomendaciones de la evaluación.
& Informe final del proyecto.
& Documentación técnica y académica.
& \LaTeX, Visual Studio Code. \\ \hline

\end{tabular}
\end{table}
\end{landscape}

\clearpage
